\section{Method}
\label{sec:method}

This section specifies the experimental design in enough detail to reconstruct
it. Where a choice was contestable we state the alternative and why it was
rejected, because several of the conclusions below depend on choices that a
reader may reasonably want to overturn.

\subsection{Corpora and label space}
\label{sec:corpora}

We use two English corpora: RAVDESS \cite{livingstone2018ryerson}
(24 speakers, 1440 utterances, 1.48\,h, mean duration 3.70\,s) and CREMA-D
\cite{cao2014crema} (91 speakers, 7442 utterances, 5.26\,h, mean duration
2.54\,s). Both are publicly available, both carry speaker identity for
every utterance, and their categorical label sets share a six-class
intersection large enough to evaluate per-class transfer. The pairing also
admits a source-size control: capping the larger training split to the smaller
lets each direction be run
at matched training-set size (Section~\ref{sec:splits}), so a reported
asymmetry between directions is not attributable to source-training-set size
alone. Both directions are evaluated.

Both are elicited from actors rather than recorded spontaneously.
Section~\ref{sec:limitations} takes up what that does and does not license.

The shared label space is the six-class intersection
$\{$angry, disgust, fear, happy, neutral, sad$\}$, in alphabetical order, which
is also the index order of every confusion matrix we report. Two mapping
decisions affect the priors and are stated because they are contestable.

\emph{RAVDESS \texttt{calm} is merged into \texttt{neutral}}, since calm exists
only in RAVDESS and cannot transfer. Merging gives neutral 288 utterances
against 192 for every other class; dropping would leave it at 96. Neither
option is balanced. The original speech subset is also not balanced across
all eight labels: neutral has 96 utterances, while each other label has 192.
\emph{RAVDESS \texttt{surprised} is excluded} as it has no CREMA-D counterpart,
removing 192 utterances and leaving 1248. CREMA-D contributes all 7442.
Table~\ref{tab:corpora} gives the resulting statistics.

\begin{table*}[!t]
  \centering
  \caption{Corpora and the mapped six-class intersection. Speakers, utterances and durations describe the raw speech subsets; the last six columns are mapped class priors. Merging RAVDESS \texttt{calm} changes the neutral prior.}
  \label{tab:corpora}
  \begin{tabular}{lrrrrrrrrrrr}
    \toprule
    corpus & spk & utts & hours & mean dur (s) & $n$ (6-class) & \texttt{angr} & \texttt{disg} & \texttt{fear} & \texttt{happ} & \texttt{neut} & \texttt{sad} \\
    \midrule
    RAVDESS & 24 & 1440 & 1.48 & 3.70 & 1248 & 0.154 & 0.154 & 0.154 & 0.154 & 0.231 & 0.154 \\
    CREMA-D & 91 & 7442 & 5.26 & 2.54 & 7442 & 0.171 & 0.171 & 0.171 & 0.171 & 0.146 & 0.171 \\
    \bottomrule
  \end{tabular}
  \\[2pt]
  \begin{minipage}{\linewidth}\footnotesize Derived from \texttt{data/manifest.csv}. RAVDESS \texttt{surprised} (192 utterances) is excluded as having no CREMA-D counterpart; RAVDESS \texttt{calm} is merged into \texttt{neutral}. Neither the original eight-label subset nor this six-class intersection is exactly balanced.\end{minipage}
\end{table*}


Class priors are close but not identical: RAVDESS carries $0.231$ neutral and
$0.154$ for each remaining class, against CREMA-D's $0.146$ and $0.171$.
The neutral mapping changes these empirical priors. It does not establish
that all corpus differences arise from the mapping. Section~\ref{sec:decomposition}
reports prior divergence and the separate label-harmonisation controls.

\paragraph{Label-harmonisation robustness protocol.} We pre-specified two
focused controls. The first sets \path{ravdess_calm_to_neutral=false},
drops the 192 RAVDESS \texttt{calm} utterances, and retains the six-class
intersection. The second also excludes \texttt{neutral} from both corpora,
forming the five-class space \texttt{five\_no\_neutral}. Each mapping has its
own label-map hash, split specification, result ledger and run identifiers. The
controls use cached HuBERT final-layer vectors, logistic regression, five
speaker-disjoint seeds, and the \texttt{none}, \texttt{zscore},
\texttt{mean\_shift} and \texttt{coral} rungs ($\varepsilon=10.0$). Source
validation selects each classifier exactly as in the main protocol. These are
label-harmonisation robustness controls, not replacements for the full
classifier grid.

\subsection{Splits and leakage control}
\label{sec:splits}

Each corpus pair is partitioned into four roles. The source corpus yields
\texttt{source\_train} and \texttt{source\_val}; the target corpus yields
\texttt{target\_adapt} and \texttt{target\_test}. Partitioning is
\textbf{speaker-disjoint} throughout: no speaker appears in more than one role.

The separation between \texttt{target\_adapt} and \texttt{target\_test} is the
central design constraint. Alignment methods require target-domain features by
construction, and they are given \texttt{target\_adapt} only.
\texttt{target\_test} is read exactly once per run, at scoring time, and never
by any fitted object. This is enforced rather than described: every alignment
object records the utterance identifiers it was fitted on, and an assertion
runs on the \emph{real fitted object in every run}, not on a mock and not once
at start-up. Model selection likewise never sees target data; the selection
routine is not passed target features at all, and the target score is computed
afterwards, by the caller, from a model already chosen. Four leakage assertions
hold over all 20 splits (4 pairs $\times$ 5 seeds): no group appears in more
than one split within a corpus; no utterance appears in both a source and a
target split; \texttt{target\_test} never appears in any fitted alignment
object; and the label mapping is a pure function.

At seed 0, RAVDESS $\rightarrow$ CREMA-D gives 988 / 260 / 3765 / 3677
utterances across the four roles, and CREMA-D $\rightarrow$ RAVDESS gives
5972 / 1470 / 624 / 624. Sizes vary slightly across seeds within the
speaker-disjoint constraint.

\paragraph{Matched-$n$ reverse direction.} The two directions differ sixfold in
training-set size, so any asymmetry between them would be confounded with the
amount of source data. We therefore cap \texttt{source\_train} to the smaller
direction's size per seed, subsampling CREMA-D from 5972 to 988: stratified by
class with largest-remainder allocation, so the prior is preserved to within
one utterance per class, and drawn from within the existing speaker-disjoint
split, so utterances are only removed and never moved between roles. Every
leakage guarantee therefore still holds, and the cap is derived per (pair, seed)
rather than hard-coded. \textbf{Both directions reported here are matched in
source-training size only}; a full-$n$ reverse arm was not run.
Source-validation, target-adaptation and target-test sizes remain different
between directions. Consequently this control removes one important
data-quantity confound but does not license an attribution of every directional
asymmetry to corpus properties.

\subsection{Features}
\label{sec:features}

Audio is resampled to 16\,kHz mono and peak-normalised. Three self-supervised
base models are used as frozen feature extractors: HuBERT Base
\cite{hsu2021hubert}, wav2vec 2.0 Base \cite{baevski2020wav2vec} and WavLM Base
\cite{chen2022wavlm}, each providing a CNN encoder output plus twelve
transformer layers, i.e.\ 13 hidden states of dimension 768.

\textbf{All 13 layers are cached}, not only the final one. Final-layer pooling
is the common default and is what the earlier version of this work used;
caching every layer makes depth a searchable condition instead.
Section~\ref{sec:results-aggregation} shows that this choice is worth more than
any choice \emph{among} alignment methods, though less than the decision to
align at all; so it is the axis practitioners may leave at its default while
tuning an axis whose observed differences are smaller in this study.

Two poolings are stored. \emph{Mean pooling} over time gives one 768-d vector
per utterance per layer, used by the four non-sequential classifiers.
\emph{Segment pooling} into 8 uniform temporal segments gives an
$8 \times 768$ sequence, so the Transformer head receives genuine sequence
structure; this costs $8\times$ storage and is the reason that arm exists at
all. Extraction runs at batch size 1 so no utterance is padded against a longer
neighbour: padding would place silence inside the pooled average, and would do
so unevenly across corpora, since RAVDESS utterances average $1.16$\,s longer
than CREMA-D's. A classical MFCC branch is also cached (13 coefficients with
deltas and delta-deltas, mean- and std-pooled, 78 dimensions) but is unused
here.

\subsection{The alignment ladder}
\label{sec:ladder}

The core comparison contains six feature-space conditions. We retain the term
``rung'' as a condition label, not as a claim that these are nested
moment-matching interventions (Section~\ref{sec:affine-factorisation}). Let
$X_s \in \mathbb{R}^{n_s \times d}$ be source-train features and $X_t$ be
\texttt{target\_adapt} features, with means $\mu_s, \mu_t$ and covariances
$C_s, C_t$. All six rungs are fitted on $X_s$ and $X_t$ only.

\begin{description}
  \item[\texttt{none}] The identity. $x \mapsto x$.
  \item[\texttt{zscore}] Per-dimension standardisation of each domain to zero
    mean and unit variance: $x \mapsto (x - \mu)\oslash\sigma$ within each
    domain. Matches the first and second marginal moments per dimension, with
    no cross-dimension terms.
  \item[\texttt{mean\_shift}] $x \mapsto x + (\mu_t - \mu_s)$. Matches the
    first moment only.
  \item[\texttt{coral}] Whitening by the source covariance and recolouring by
    the target's \cite{sun2016deep}:
    $x \mapsto (x - \mu_s) C_s^{-1/2} C_t^{1/2} + \mu_t$. Matches first and
    second moments including cross-terms in the unregularised idealisation.
    With shrinkage, the identity is $M^\top\tilde C_sM=\tilde C_t$;
    empirical covariances need not match exactly.
  \item[\texttt{mkmmd\_diag}] A learned diagonal affine map minimising a
    multi-kernel MMD objective, defined below.
  \item[\texttt{mkmmd\_full}] The same objective with a full matrix $W$.
\end{description}

\paragraph{Mean matching.} For a linear kernel, the squared distance between
empirical kernel means is $\lVert\mu_s-\mu_t\rVert^2$. Within translations,
its minimising offset is $\mu_t-\mu_s$. Among unrestricted affine maps this
does not uniquely select a translation: any matrix $A$ can match the means
with offset $\mu_t-\mu_sA$. The empirical-mean expression is not the unbiased
finite-sample estimator used for the RBF diagnostics. We name the operation
\texttt{mean\_shift}, not ``MMD''.

\paragraph{Regularisation of CORAL.} With $d = 768$ and roughly one thousand
source-train samples, the rank bound $\min(d,n-1)$ does not itself imply
singularity. Shrinkage stabilises whitening when small eigenvalues make it
ill-conditioned. It is
scale-aware, $\tilde{C} = C + \varepsilon\,\frac{\mathrm{tr}(C)}{d} I$:
anchoring to the mean eigenvalue matters because a fixed absolute ridge means
something different for every backbone and layer. $\varepsilon$ is a
\emph{searched} axis over $\{10^{-4}, \ldots, 10\}$ in decades, selected on
\texttt{source\_val} and recorded on every row, with a Ledoit--Wolf variant
\cite{ledoit2004covariance} as a
parameter-free alternative. Condition number and effective rank are recorded
after regularisation; a matrix still singular there raises rather than being
pseudo-inverted.

\paragraph{Multi-kernel MMD.} The two MK-MMD rungs learn an affine map
$(W, b)$ minimising
\begin{equation}
  \mathcal{L}(W, b) \;=\;
  \mathrm{MMD}^2_k\!\left(X_s W^\top + b,\; X_t\right)
  \;+\; \lambda \lVert W - I \rVert_F^2 ,
  \label{eq:mkmmd}
\end{equation}
where the implementation stores $W$ as the transpose of the row-vector map
used in Section~\ref{sec:affine-factorisation}, and $k$ is a sum of Gaussian
kernels at bandwidth multipliers
$\{0.25, 0.5, 1, 2, 4\}$ of the median pairwise distance
\cite{gretton2012kernel}, and $\mathrm{MMD}^2$ is the unbiased estimator.
\texttt{mkmmd\_diag} constrains $W$ to be diagonal; \texttt{mkmmd\_full} does
not. The identity anchor regularises a full matrix with $\approx 590$k
parameters fitted on about a thousand source samples. It penalises $W-I$,
not $b$. Even an ideal large-$\lambda$ limit therefore leaves a translation
to optimise; it does not necessarily recover \texttt{none}. No asymptotic
optimisation guarantee is claimed for the finite-step procedure.
$\lambda$ is searched over
$\{10^{-3}, \ldots, 10^{2}\}$ in decades.

Optimisation uses Adam \cite{kingma2015adam} for 200 steps at batch size 256,
warm-started from the
CORAL solution (full) or a diagonal moment match (diagonal). The learning rate
is $\eta = \eta_0 / \sqrt{P}$ with $\eta_0 = 0.01$ and $P$ the parameter count:
Adam normalises per parameter, so a fixed rate gives a step whose Frobenius
norm grows as $\sqrt{P}$; for a full $W$, norm $\approx 0.77$ per iteration,
which does not converge.

\paragraph{Fallback, and why it must be reported.} An unpenalised MMD check
on the first 400 observations compares the candidate with its warm start.
If the candidate worsens this check, the returned map reverts to the
warm start. This differs from the penalised minibatch training objective;
rejection is not proof of failure to improve that objective. Optimiser
diagnostics describe the attempted fit, while target scores use the returned
map. Across the main grid the check rejects the candidate
in \textbf{64.9\%} of \texttt{mkmmd\_diag} runs and \textbf{55.2\%} of
\texttt{mkmmd\_full} runs. In those cells the fitted map is the warm start:
the diagonal moment match for \texttt{mkmmd\_diag}, and CORAL for
\texttt{mkmmd\_full}. Reporting either as a distinct learned map without saying
so would overstate what was actually fitted. The rates accompany the ladder
result in Section~\ref{sec:results-ladder} and the generated results document.
They qualify every MK-MMD comparison in this paper: the experiment evaluates
these affine maps, this warm start, and this optimisation budget. It does not
establish what a differently optimised MK-MMD map would achieve.

\subsection{Measuring distribution discrepancy}
\label{sec:discrepancy}

Alignment methods are usually justified by the discrepancy they remove, so the
discrepancy statistic needs to be specified as carefully as the methods.

We report the raw unbiased $\mathrm{MMD}^2$ between aligned source and target
and its effect size, obtained by dividing by the mean absolute $\mathrm{MMD}^2$
over five random source half-splits. The unbiased estimator has null mean zero,
so its signed mean is not a usable denominator. The effect size aids
interpretation of finite-sample variation; it does not replace the raw
statistic.

\paragraph{A common reference geometry.} For the ladder diagnostic, one ZCA
basis is derived from unaligned \texttt{source\_train} features with shrinkage
$\varepsilon = 10^{-2}$. One median RBF bandwidth is then derived from those
same source-train features in that basis. Both quantities are fixed before an
alignment map is fitted and applied identically to every rung for that pair,
seed, backbone and layer aggregation. This makes the raw trajectories a
comparison under one stated kernel and coordinate system. The shrinkage avoids
amplifying near-null directions in the source covariance. A
scale-changing map can nonetheless drive a fixed-bandwidth RBF kernel into
saturation. We report that failure mode rather than treating a near-zero
statistic as evidence of distributional equality.
Each MMD evaluation uses a seeded cap of 128 utterances per sample set; the cap
is identical for the marginal, class-conditional and null computations and is
recorded in every diagnostic row.

\paragraph{Adaptive geometry as a sensitivity analysis.} We separately report
an own-geometry calculation that derives the median bandwidth from each
compared feature pair. The 13-layer frame-sensitivity sweep also holds the ZCA
basis fixed while recomputing that adaptive bandwidth; we call it a
\emph{reference-basis} calculation, not a fully fixed reference geometry. The
two procedures answer different measurement questions, so neither is treated
as the unique definition of distributional closeness.

\paragraph{Proposition: median-heuristic scale invariance.} Let
$\hat\sigma(X,Y)$ be the median pairwise Euclidean distance used by the RBF
kernel. For a nonzero scalar $a$,
$\hat\sigma(aX,aY)=|a|\hat\sigma(X,Y)$ and therefore
\begin{equation}
  k_{|a|\hat\sigma}(ax,ay)=k_{\hat\sigma}(x,y).
  \label{eq:median-scale-invariance}
\end{equation}
Each pairwise distance gains the factor $|a|$, which proves the first identity;
substitution in the RBF exponent proves the second. Thus an adaptive
median-heuristic MMD is invariant to a global rescaling applied to both samples.
This property does not extend to an arbitrary affine map or to a map applied to
only one sample.

\paragraph{Class-conditional diagnostic.} For each class $k$, the analysis
computes the same unbiased multi-kernel MMD$^2$ between aligned
$X_s \mid y=k$ and $X_t \mid y=k$ in the common reference geometry. A class is
undefined rather than estimated when either side has fewer than 50 examples.
The report also gives a class-specific half-split effect size. That normaliser
is not the marginal normaliser, so we do not form or interpret a
conditional-to-marginal ratio. The diagnostic is not a probability
decomposition or an estimate of $P(y \mid x)$.

\subsection{Classifiers and the search budget}
\label{sec:classifiers}

Five families are evaluated: logistic regression, linear SVM, RBF SVM, an MLP,
and a small Transformer head over the 8-segment sequence. Each is searched with
\textbf{exactly the same number of trials}; 20 random-search draws; and
the trial count is recorded on every result row. This controls trial count,
not effective model capacity or coverage of equally favourable settings.

Hyperparameter candidates are drawn by seeded random search. A failed fit
receives score $-\infty$ and is not replaced by an additional draw.

Two further constraints. First, \textbf{no standardisation inside any
classifier}. This preserves the frozen raw-coordinate baseline. A scaler
fitted on source training and shared across domains would be a different
control; it would not equal domain-wise \texttt{zscore}. Consequently this
grid does not isolate alignment from classifier conditioning. Second, the
iteration cap for the iterative solvers is a fixed convergence budget
($2\times10^4$) rather than a searched hyperparameter, as it was in an earlier
revision: searching it lets a trial win selection by having stopped early
rather than by fitting better. Convergence is asserted rather than warned
about; a fit that reaches the cap is recorded as a failed trial and cannot
be selected; and across the whole grid \textbf{0 of 99{,}720 trials failed
to converge}.

Layer aggregation is a searched condition with three settings: \texttt{last}
(the common default), \texttt{layer:6} (a fixed mid-stack layer), and
\texttt{weighted} (a softmax over all 13 states learned jointly with the head).
The last is implemented only for the two torch families, whose training loops
jointly optimise the layer weights and head. It is not searched for the
scikit-learn classifiers in this implementation.

The Transformer arm is a deliberately \textbf{reduced-seed arm}: two seeds, the
primary direction only, one inner-grid setting per rung. Its per-cell cost is
15--40$\times$ the other families', and carrying it at full width would have
spent more compute on it than on every other axis combined. The full-grid
headline and ladder summaries include its available cells, so the forward
candidate set is not identical across seeds. The matched-direction comparison
in Section~\ref{sec:results-direction} excludes it. The translation audit
uses only the RBF-SVM and logistic-regression families and is unaffected.

\subsection{Selection protocol}
\label{sec:selection}

We report two protocols side by side, and the distinction between them is one
of the paper's results rather than a presentational choice. \textbf{Validated}
is the configuration with the highest \texttt{source\_val} macro-F1 within a
given (pair, seed), scored on \texttt{target\_test}; it is what a practitioner
could actually obtain, using no target labels at any point. \textbf{Oracle} is
the highest \texttt{target\_test} macro-F1 present in the grid for that
(pair, seed): a \textbf{post-hoc target-selected benchmark unavailable to the
deployable source-only selection protocol}. The
validated-to-oracle difference is called the \emph{oracle gap}. It does not
identify a causal effect of source-side selection because it also contains
finite-target-sample noise, search-space size, and optimism from maximising the
target score.

The gap between them is the quantity of interest. Reporting only the oracle
amounts to selecting on target scores. The earlier version of this work did
this. It makes a method look as good as its best run.

\paragraph{How the ladder handles unequal inner grids.} Each result row already
contains the best of its classifier family's 20 source-validation trials. For a
ladder summary, we then hold (pair, seed, backbone, layer aggregation,
classifier family) fixed and select that rung's best remaining inner setting
($\varepsilon$ for CORAL or $\lambda$ for MK-MMD) by
\texttt{source\_val} macro-F1. Target predictions from only these selected
cells enter the paired cluster bootstrap. The row counts in Table~\ref{tab:ladder}
are therefore candidate rows before this second selection, not equally weighted
observations in a rung mean; they differ because the rungs legitimately have
different inner grids.

\subsection{Metrics, baselines and inference}
\label{sec:metrics}

The primary metric is macro-F1 on \texttt{target\_test}. Accuracy and
unweighted average recall are also recorded on every row.

\textbf{Every performance number is reported against its baselines.} An analytic
chance baseline (the expected macro-F1 of a uniformly random classifier, given
the realised target-test priors) and a majority-class baseline are computed on
every row from the realised distribution rather than by simulation. Here the
chance baseline is $0.1665$ (RAVDESS $\rightarrow$ CREMA-D) and $0.1656$
(CREMA-D $\rightarrow$ RAVDESS), and the majority baseline is $0.0425$ and
$0.0444$. A macro-F1 without its baselines is not interpretable, and one number
below its chance baseline is below the expected random-classifier score. Because
the baseline is pair-dependent,
\textbf{no figure in this paper averages macro-F1 across pairs}.
Table~\ref{tab:floors} collects the split sizes and both baselines.

\begin{table}[tb]
  \centering
  \caption{Split sizes and the floors every macro-F1 in this paper is read against. Both directions are matched-$n$: CREMA-D source-train is subsampled from 5972 to match RAVDESS, so a reported asymmetry cannot be a training-set size effect.}
  \label{tab:floors}
  \begin{tabular}{lrrrrrr}
    \toprule
    pair & source\_train & source\_val & target\_adapt & target\_test & chance & majority \\
    \midrule
    RAVDESS $\rightarrow$ CREMA-D & 988 & 260--260 & 3752--3765 & 3677--3690 & 0.1665 & 0.0425 \\
    CREMA-D $\rightarrow$ RAVDESS & 988 & 1470--1476 & 624--624 & 624 & 0.1656 & 0.0444 \\
    \bottomrule
  \end{tabular}
  \\[2pt]
  \begin{minipage}{\linewidth}\footnotesize Ranges span the five seeds within the speaker-disjoint constraint. Floors are analytic from the realised \texttt{target\_test} priors, not simulated. Because the chance floor is pair-dependent, no result in this paper averages macro-F1 across pairs.\end{minipage}
\end{table}


\paragraph{Intervals.} Target-score intervals come from a \textbf{paired
cluster bootstrap} with 2000 replicates \cite{cameron2008cluster}. In each
replicate, we draw the five seed identifiers with replacement. For each drawn
seed, we draw that seed's target-test speaker identifiers with replacement,
keeping every utterance from each selected speaker. Both comparison arms use
the same seed and speaker draws, then their macro-F1 difference is formed. This
resamples the split, search and initialisation variation carried by seeds and
the clustered target-test sampling variation carried by speakers. It never
resamples individual utterances independently.

\paragraph{Why the source-selected summary uses a different interval.} The
paired cluster bootstrap requires a fixed arm: a condition that can be
recomputed on any resampled set of speakers and seeds. The source-selected
summary in Table~\ref{tab:headline} is not such an arm. Its value for a seed is
the target macro-F1 of \emph{whichever} configuration source-side validation
chose on that seed, and the chosen configuration differs across seeds; in the
forward direction the protocol selects \texttt{none} on two of five seeds and an
aligned rung on the other three (Section~\ref{sec:results-selection}). Each seed
therefore contributes one number describing a complete deployment decision, and
those five numbers are summarised by a mean and a 95\% $t$-interval over seeds.

That interval estimates the between-seed variability of the deployable
selection outcome, which includes both the split realisation and the change of
selected configuration it induces. It does \textbf{not} incorporate any further
within-seed resampling of target-test speakers, so it is not comparable to the
paired intervals elsewhere in this paper and is not offered as a better
estimator. It answers a different question: how much the end-to-end protocol
varies from one split to another. Its width is also limited by having five
observations. Fixed-arm contrasts, where both conditions can be recomputed on
identical seed and speaker draws, use the paired cluster bootstrap throughout.
The same estimator is applied to the oracle column and to their difference, so
the three columns of that table are on a common basis.

\paragraph{What the oracle column is.} The oracle is the best target macro-F1
present in the grid for that pair and seed. It is selected using target-test
scores and is therefore not attainable by any source-only protocol; it is a
benchmark for the grid as configured, not a result. Its interval carries the
same maximisation: the maximum of a set of noisy scores is biased upward, so
the oracle column is an optimistic summary of this target set rather than an
estimate of achievable performance. The difference between the two columns is
reported as an \emph{oracle gap} and is descriptive. It is not the causal
effect of using source-side selection, because it also contains target-set
maximisation, finite-sample variation and the extent of the search space.

The discrepancy columns are properties of a fitted map, not of a prediction.
The ladder discrepancy columns are candidate-row means without intervals.
Other discrepancy summaries state their own aggregation. The frame sweep's
cell-and-seed correlations are descriptive because cells share seeds and
feature sets; pooled independence-based intervals are not used as evidence
for a population-level sign change.

\paragraph{AI-assisted code development.} OpenAI Codex was used to assist with
code scaffolding and maintenance. The authors specified the experimental
protocol, reviewed and edited all generated code, ran the tests and experiments,
and verified every reported result against the generated artifacts.

\paragraph{Multiple comparisons.} A primary comparison family of seven
contrasts $\times$ two directions was \textbf{pre-specified in code before the
frozen confirmatory grid was run}: the five ladder rungs against \texttt{none}, and
two layer-aggregation contrasts. Holm-Bonferroni correction
\cite{holm1979sequential} is applied across exactly those 14 tests.
Holm rather than Bonferroni because it is uniformly more powerful at the same
family-wise error rate, and rather than a false-discovery-rate procedure
because the family is small and pre-specified, where controlling the
family-wise rate is the stricter guarantee.

Comparisons outside that family are reported as such, and are never presented
as though they had been pre-specified. Where such a comparison supports a
\emph{null} claim; that some difference is negligible; we report a bound
rather than a test, since failing to reject a null is not evidence for it.
Where one supports a positive claim, correction is applied within that
comparison's own set and both the raw and corrected $p$ are given, so that a
result surviving only before correction is visible as such.

\subsection{Experimental scope}
\label{sec:scope}

The main grid is a reduced factorial; 2 directions $\times$ 3 backbones
$\times$ 5 seeds $\times$ 6 rungs $\times$ layer aggregation $\times$ 5
classifier families, with inner grids over $\varepsilon$ and $\lambda$. This
totals 4986 runs. A smoke gate and a screening stage pruned
only the inner grids on \texttt{source\_val} evidence alone. Three probes
supplement it: a 13-layer sweep (2340 runs), a CORAL shrinkage asymptote probe
(120 runs) and a shift decomposition (120 records). No run failed. Results here
carry the frozen configuration tag \texttt{grid-freeze-v3};
Section~\ref{sec:reproducibility} describes the provenance machinery.
